\documentclass{leveldoc}
\DocUses{foundations,slingshot,first_coil}
\DocFingerprint{ace70898a15db445}

\begin{document}

\section{What the level asks}\label{sec:asks}

The game's text:
\leveldescription

Three positive ions start together at $(\val{xZero}, \val{yZero})$ along $+x$ with the kinetic
energies $\val{TSlow}$, $\val{TMid}$ and $\val{TFast}$. A rectangular coil around the board,
from $(\val{cx0}, \val{cy0})$ to $(\val{cx1}, \val{cy1})$, fills it with a magnetic field
perpendicular to the board, which alone would bend every ion down into the coil's lower side.
The player places up to four charges from the listed magnitudes in the region
$\val{regX0} \le x \le \val{regX1}$, $\val{regY0} \le y \le \val{regY1}$, to supply an electric
field. The middle ion must enter its detector, $\val{xDetLoMid} \le x \le \val{xDetHiMid}$,
$\val{yDetLoMid} \le y \le \val{yDetHiMid}$, moving within $\pm\val{coneHalf}$ degrees of $+x$;
the slow ion the detector above it, $\val{yDetLoSlow} \le y \le \val{yDetHiSlow}$, and the fast
ion the one below, $\val{yDetLoFast} \le y \le \val{yDetHiFast}$ (both at the same $x$, without
a direction window). A flight ends when the ion touches a charge or the coil's wire, leaves the
world, or enters its detector, which a direction outside the window turns from an arrival into
a rejection (\xsecref{foundations}{sec:board}, \physref{6.1}).

The reference places a negative charge above the path near its start and a positive one below
it further on: both push the ions upwards, against the magnetic force. The electric force on an
ion does not depend on its speed; the magnetic force grows with it. Where they balance for the
middle ion, $qE = qvB$, it goes straight; the slow ion is lifted, the fast one pushed down. This
document derives the rectangular coil's field from the Biot--Savart law and integrates the three
flights numerically.

\section{Setup}\label{sec:setup}

The level file, field by field (the appendix reproduces it):
\setuptable

In symbols: each ion has charge $q = \val{q}$, mass $m = \val{m}$ and radius $0$; $c = \val{c}$.
Their speeds are $v = \val{vSlow}$, $\val{vMid}$ and $\val{vFast}$ (\xeqref{foundations}{eq:launch}).
The coil is the rectangle's four straight sides, carrying the strength $\kappa = \val{kappa}$
(\physref{2.2}) counter-clockwise; its wire has radius $\val{wire}$. The reference charges are
$Q_1 = \val{QOne}$ at $(\val{xOne}, \val{yOne})$ and $Q_2 = \val{QTwo}$ at $(\val{xTwo}, \val{yTwo})$,
radius $R = \val{R}$. The world is $-1 \le x \le \val{xWorld}$, $-1 \le y \le \val{yWorld}$.

\begin{figure}[htbp]
\centering
\begin{tikzpicture}
\begin{axis}[axis equal image, width=0.92\linewidth, xmin=-1.6, xmax=31.6, ymin=-1.6, ymax=21.6,
  xtick={0,5,...,30}, ytick={0,5,...,20}, tick label style={font=\small},
  xlabel={$x$}, ylabel={$y$}, label style={font=\small}]
  \draw[black!8] (0,0) grid (30,20);
  \draw[black!45, dashed] (-1,-1) rectangle (\rawval{xWorld},\rawval{yWorld});
  \draw[black!70, thick] (\rawval{cx0},\rawval{cy0}) rectangle (\rawval{cx1},\rawval{cy1});
  \draw[black!30, dotted] (\rawval{regX0},\rawval{regY0}) rectangle (\rawval{regX1},\rawval{regY1});
  \fill[green!50!black, fill opacity=0.18] (\rawval{xDetLoSlow},\rawval{yDetLoSlow}) rectangle (\rawval{xDetHiSlow},\rawval{yDetHiSlow});
  \draw[green!40!black] (\rawval{xDetLoSlow},\rawval{yDetLoSlow}) rectangle (\rawval{xDetHiSlow},\rawval{yDetHiSlow});
  \fill[green!50!black, fill opacity=0.18] (\rawval{xDetLoMid},\rawval{yDetLoMid}) rectangle (\rawval{xDetHiMid},\rawval{yDetHiMid});
  \draw[green!40!black] (\rawval{xDetLoMid},\rawval{yDetLoMid}) rectangle (\rawval{xDetHiMid},\rawval{yDetHiMid});
  \fill[green!50!black, fill opacity=0.18] (\rawval{xDetLoFast},\rawval{yDetLoFast}) rectangle (\rawval{xDetHiFast},\rawval{yDetHiFast});
  \draw[green!40!black] (\rawval{xDetLoFast},\rawval{yDetLoFast}) rectangle (\rawval{xDetHiFast},\rawval{yDetHiFast});
  \addplot[thick, blue!70!black] table {pathslow.dat};
  \addplot[thick, green!40!black] table {pathmid.dat};
  \addplot[thick, red!70!black] table {pathfast.dat};
  \fill[blue!75!black] (\rawval{xOne},\rawval{yOne}) circle[radius=\rawval{R}];
  \fill[red!75!black] (\rawval{xTwo},\rawval{yTwo}) circle[radius=\rawval{R}];
  \fill (\rawval{xZero},\rawval{yZero}) circle[radius=0.15];
\end{axis}
\end{tikzpicture}
\caption{The rectangular coil (black), the reference's negative charge above the path (blue)
and positive charge below it (red), and the three flights: slow (blue) lifted into the upper
detector, middle (green) straight into the middle one, fast (red) pushed into the lower one.}
\label{fig:board}
\end{figure}

\section{Principles and model}\label{sec:model}

\begin{principles}
  \item The Lorentz force $q(\vect{E} + \vect{v}\times\vect{B})$ in the relativistic equation of
    motion \xeqref{foundations}{eq:eom}; the point charges' field \xeqref{foundations}{eq:coulomb};
    the energy \xeqref{foundations}{eq:energy}.
  \item The Biot--Savart law for a closed circuit (\Jackson{§5.2}) in the game's units
    (\physref{2.2}).
  \item The board's rules and the detector's direction window (\xsecref{foundations}{sec:board},
    \physref{6.1}); the classical Runge--Kutta method, cubic Hermite interpolation and the error
    estimate by halving the step (\xsecref{foundations}{sec:numerics}), as in
    \xsecref{slingshot}{sec:method}.
\end{principles}

The assumptions are First coil's (\xsecref{first_coil}{sec:model}): the coil and the charges are
held in place, the wires are thin lines, the ions are points, and their radiation is not taken
from them.

\section{Derivation}\label{sec:derivation}

\subsection{A straight side's field}\label{sec:segment}

Let a side run from $\vect{a}$ to $\vect{b}$ along the unit vector $\uvec{t}$, and let the field
point $\vect{x}$ lie in the plane at the distance $d$ from the side's line, its foot at the
origin of the arc length $s$ along $\uvec{t}$; the side spans $s_a \le s \le s_b$. An element
$\dd\vect{l} = \uvec{t}\,\dd s$ at $s$ contributes, by the Biot--Savart law,
$\kappa\,\dd\vect{l}\times(\vect{x} - \vect{l})/\abs{\vect{x} - \vect{l}}^3$, whose size is
$\kappa\,d\,\dd s/(d^2 + s^2)^{3/2}$ and whose direction is $\pm\uvec{z}$, the sign of
$(\uvec{t}\times(\vect{x} - \vect{a}))_z$. With
$\int \dd s/(d^2 + s^2)^{3/2} = s/(d^2\sqrt{d^2 + s^2})$,
\begin{equation}\label{eq:segment}
  B_z = \pm\frac{\kappa}{d}\left(\frac{s_b}{\sqrt{d^2 + s_b^2}} - \frac{s_a}{\sqrt{d^2 + s_a^2}}\right).
\end{equation}
The coil's field is the sum over its four sides. A direct Simpson quadrature of the Biot--Savart
integral along each side agrees with \eqref{eq:segment} to a relative $\val{segAgree}$. For an
infinitely long side $s_a \to -\infty$, $s_b \to \infty$, and \eqref{eq:segment} becomes
$2\kappa/d$, the long straight wire's field: at the rectangle's centre the two long sides,
$\val{dSide}$ away, would give $\val{Bstrip}$; with the short sides the field there is
$\val{Bmid}$.

\subsection{The balance}\label{sec:balance}

For an ion moving along $+x$ in $B_z > 0$ the magnetic force $q\vect{v}\times\vect{B} =
qB_z(v_y, -v_x)$ points down, $-qvB_z$; an electric field $E_y$ adds $qE_y$. They cancel where
\begin{equation}\label{eq:balance}
  v = \frac{E_y}{B_z},
\end{equation}
independently of the mass and the charge: a Wien filter selects a speed. A slower ion feels the
electric force more and rises, a faster one sinks. With point charges $E_y/B_z$ is not constant
along the path (the table, on the launch line, with the reference charges); the middle ion goes
straight on average, not everywhere, and the detector's window of $\pm\val{coneHalf}^\circ$ asks
for the net bend to cancel:
\begin{center}\small\val{axis}\end{center}
Here $\vect{E}$ comes from the charges alone; the coil's current is static and makes none.

\subsection{The method}\label{sec:method}

The three flights are integrated by \xeqref{foundations}{eq:rk4} at $h = \val{H}$, $B_z$ from
\eqref{eq:segment} at every evaluation, with the events located on the cubic Hermite
interpolant: a charge's disc, the coil's wire (within $\val{wire}$ of a side), the world's edge,
the detector, where the arrival direction is tested against the window. The energy
\xeqref{foundations}{eq:energy} is the check; the largest step-halving estimate of an end state
is $\val{errMax}$.

\section{The reference solution by hand}\label{sec:hand}

\begin{procedure}
  \item\label{st:fields} \emph{The fields.} $B_z$ by \eqref{eq:segment} summed over the four
    sides; $\vect{E}$ from the two charges.
  \item\label{st:slow} \emph{The slow ion} ($v = \val{vSlow}$, below the balance) is lifted and
    enters its detector through the \val{faceSlow} face at $(\val{xSlow}, \val{ySlow})$ at
    $t = \val{tSlow}$, moving $\val{angSlow}$ degrees above $+x$.
  \item\label{st:mid} \emph{The middle ion} ($v = \val{vMid}$) enters its detector through the
    \val{faceMid} face at $(\val{xMid}, \val{yMid})$ at $t = \val{tMid}$, moving $\val{angMid}$
    degrees from $+x$: inside the window, so it arrives.
  \item\label{st:fast} \emph{The fast ion} ($v = \val{vFast}$) is pushed down and enters its
    detector through the \val{faceFast} face at $(\val{xFast}, \val{yFast})$ at $t = \val{tFast}$,
    moving $\val{angFast}$ degrees from $+x$.
  \item\label{st:checks} \emph{The checks.} The energy drifts by at most $\val{wDevFast}$; the
    step estimates are at most $\val{errMax}$.
\end{procedure}

\begin{result}[The reference flights]
The three ions enter their detectors at $(\val{xSlow}, \val{ySlow})$, $(\val{xMid}, \val{yMid})$
and $(\val{xFast}, \val{yFast})$, the middle one within its window. The level is solved.
\end{result}

\section{Comparison with the game}\label{sec:comparison}

\comparisontable

All differences are a few $10^{-11}$ or less, the game's verified accuracy (Foundations §6);
our step estimates are below $10^{-12}$. The tolerance, $10^{-9}$, is ten times the larger of the
two computations' estimates (or the game's own preview--verification gap where larger), fixed
before comparing. The middle ion's direction, $\val{angMid}$ degrees, is far inside the window,
so its outcome is not marginal.

\section{Other solutions}\label{sec:others}

On the empty board, the second placement in \code{placements.json}, the magnetic force bends all
three ions down into the coil's lower side: the middle one strikes the wire at
$(\val{bareXMid}, \val{bareYMid})$ (\val{bareOutMid}; the table has all three). A solution must
supply an upward electric field that balances the middle ion's magnetic force over the whole
path: one charge's field, strong near it and weak elsewhere, can sort three speeds but rarely
lets the middle one leave straight, which is why the level asks for the direction. Charges of
either sign above and below the path can share the work; a map of solutions would take the
game's verdicts, which this document does not reproduce.

\section{What the game computes}\label{sec:engine}

The game evaluates a polygonal coil's field as the sum of its straight segments' exact fields
(\physref{2.2}), of which \eqref{eq:segment} is the in-plane case, and integrates the equations of
motion with its adaptive eighth-order Runge--Kutta method (\physref{5}, \physref{6},
\physref{7}), testing the arrival direction against the window (\physref{6.1}); the differences
in the comparison are the two integrations' errors.

\section{Further reading}\label{sec:reading}

The Biot--Savart law and a straight wire's field: \Jackson{§5.2, §5.3}; motion in crossed
uniform fields and the drift $\vect{E}\times\vect{B}/B^2$: \Jackson{§12.3}. Magnetic fields and
cyclotron motion: \levelref{first_coil}.

\levelfileappendix

\end{document}
